% SPDX-License-Identifier: Apache-2.0
% covers: NoBeats
\datasheet{NoBeats}{The write-beat channel of a host whose client only
reads, held and never offered on.}

\dsfacts{\code{lib/parts/src/dma.rs}}{\code{txhdl\_parts::dma::NoBeats}}%
{\code{//docs:datasheets}}

\subsection*{Function}

A host takes write beats from its client whether the client writes or
not, and a board joins every channel it makes to a unit: an end that
nothing holds is refused rather than left floating. A fetch engine
never writes, so something has to hold its host's write-beat channel.
\code{NoBeats} holds it and offers nothing on it.

\subsection*{Parameters}

None. The beat is the tree's write beat, \code{W<32, 4>}.

\dstables{NoBeats}

\subsection*{Behaviour}

\begin{itemize}
\item It never offers a beat. In the netlist its valid is low for
ever, which is what a client that never writes would drive.
\item \code{idle} is its one register, and nothing sets it. It is
there only because a unit is a struct of registers; it keeps nothing.
\end{itemize}

\subsection*{Verification}

It has no example and no test of its own; what it does is nothing. It
is in the board's netlist, and the board tests in
\code{cpu/vreteno/tests/board.rs} run with it in place.

\subsection*{Used by}

\code{Board}, as \code{fnobeats}, on the write-beat channel of the
fetch engine's host \code{fhost}.

\subsection*{Limits}

A host whose client does write must not be given one: its writes would
wait for beats that never come.

The mirror for a client that only writes is \code{NoReads}.
